\chapter{Phụ lục tra cứu kỹ thuật}
\label{ch:phu-luc}

\section{Danh mục 34 biến số trong không gian đặc trưng $F_4$}
\label{sec:app-feature-inventory}

Bảng \ref{tab:feature-inventory} trình bày đầy đủ danh mục 34 biến số đầu vào trong không gian đặc trưng $F_4$, bao gồm mã định danh của CDC NHANES, tên biến, kiểu dữ liệu và mô tả lâm sàng tương ứng.

\begingroup
\setlength{\tabcolsep}{3pt}
\begin{longtable}{>{\raggedright\arraybackslash}p{2.9cm} >{\raggedright\arraybackslash}p{2.6cm} >{\raggedright\arraybackslash}p{2.5cm} >{\raggedright\arraybackslash}p{6.3cm}}
\caption{Danh mục hệ thống 34 đặc trưng trong không gian $F_4$ theo giao thức chính thức} \label{tab:feature-inventory} \\
\toprule
\textbf{Mã biến} & \textbf{Nhóm đặc trưng} & \textbf{Kiểu dữ liệu} & \textbf{Mô tả lâm sàng / ý nghĩa} \\
\midrule
\endfirsthead
\multicolumn{4}{c}{\small\textbf{Bảng \thetable\ (tiếp theo). Danh mục hệ thống 34 đặc trưng trong không gian $F_4$}} \\
\addlinespace[0.2cm]
\toprule
\textbf{Mã biến} & \textbf{Nhóm đặc trưng} & \textbf{Kiểu dữ liệu} & \textbf{Mô tả lâm sàng / ý nghĩa} \\
\midrule
\endhead
\bottomrule
\endfoot
\multicolumn{4}{l}{\textbf{Nhóm $F_1$ (10 đặc trưng): Dấu hiệu sinh tồn, Nhân trắc \& Xét nghiệm máu}} \\
BMXBMI & Thể chất ($F_1$) & Liên tục (kg/m$^2$) & Chỉ số khối cơ thể (Body Mass Index) \\
BMXWAIST & Thể chất ($F_1$) & Liên tục (cm) & Chu vi vòng eo (đo lường béo phì trung tâm) \\
BPXPLS & Sinh tồn ($F_1$) & Rời rạc (nhịp/phút) & Tần số mạch đập khi nghỉ ngơi (60 giây) \\
SBP\_mean & Sinh tồn ($F_1$) & Liên tục (mmHg) & Huyết áp tâm thu trung bình qua các lần đo \\
DBP\_mean & Sinh tồn ($F_1$) & Liên tục (mmHg) & Huyết áp tâm trương trung bình qua các lần đo \\
LBXGH & Cận lâm sàng ($F_1$) & Liên tục (\%) & Glycohemoglobin (HbA1c - chuyển hóa đường) \\
LBXTC & Cận lâm sàng ($F_1$) & Liên tục (mg/dL) & Cholesterol toàn phần huyết thanh \\
LBDHDD & Cận lâm sàng ($F_1$) & Liên tục (mg/dL) & HDL-Cholesterol (cholesterol tỷ trọng cao) \\
LBXSCR & Cận lâm sàng ($F_1$) & Liên tục (mg/dL) & Creatinine huyết thanh (chức năng lọc của thận) \\
LBXSAL & Cận lâm sàng ($F_1$) & Liên tục (g/dL) & Albumin huyết thanh (chức năng gan, dinh dưỡng) \\
\midrule
\multicolumn{4}{l}{\textbf{Nhóm $F_2$ (19 đặc trưng): $F_1$ + 9 tiền sử bệnh lý lâm sàng}} \\
BPQ020 & Tiền sử ($F_2$) & Phân loại nhị phân & Bác sĩ từng chẩn đoán Tăng huyết áp \\
DIQ010 & Tiền sử ($F_2$) & Phân loại nhị phân & Bác sĩ từng chẩn đoán Đái tháo đường \\
MCQ160B & Tiền sử ($F_2$) & Phân loại nhị phân & Bác sĩ từng chẩn đoán Suy tim sung huyết \\
MCQ160C & Tiền sử ($F_2$) & Phân loại nhị phân & Bác sĩ từng chẩn đoán Bệnh tim mạch vành \\
MCQ160D & Tiền sử ($F_2$) & Phân loại nhị phân & Bác sĩ từng chẩn đoán Đau thắt ngực \\
MCQ160E & Tiền sử ($F_2$) & Phân loại nhị phân & Tiền sử từng bị Cơn đau tim / Nhồi máu cơ tim \\
MCQ160F & Tiền sử ($F_2$) & Phân loại nhị phân & Tiền sử từng bị Đột quỵ / Tai biến mạch máu não \\
MCQ220 & Tiền sử ($F_2$) & Phân loại nhị phân & Tiền sử từng được chẩn đoán mắc Ung thư \\
HUQ010 & Lâm sàng ($F_2$) & Thứ bậc (1--5) & Tự đánh giá tình trạng sức khỏe tổng quát \\
\midrule
\multicolumn{4}{l}{\textbf{Nhóm $F_3$ (29 đặc trưng): $F_2$ + 10 hành vi lối sống \& giấc ngủ}} \\
PAQ605, PAQ620 & Lối sống ($F_3$) & Phân loại nhị phân & Hoạt động thể lực cường độ mạnh / vừa tại nơi làm việc \\
PAQ650, PAQ665 & Lối sống ($F_3$) & Phân loại nhị phân & Hoạt động thể lực cường độ mạnh / vừa trong giải trí \\
PAD680 & Lối sống ($F_3$) & Liên tục (phút/ngày) & Thời gian ngồi tĩnh tại trung bình mỗi ngày \\
smoking\_status & Lối sống ($F_3$) & Phân loại 3 mức & Tình trạng hút thuốc lá (Never, Former, Current) \\
ALQ101 & Lối sống ($F_3$) & Phân loại nhị phân & Uống ít nhất 12 đồ uống có cồn trong năm qua \\
alcohol\_status & Lối sống ($F_3$) & Phân loại 3 mức & Tình trạng dùng rượu bia (Non, Moderate, Heavy) \\
alcohol\_qty & Lối sống ($F_3$) & Thứ bậc lượng hóa & Mức độ tiêu thụ rượu bia lượng hóa \\
SLD010H & Lối sống ($F_3$) & Liên tục (giờ) & Thời lượng giấc ngủ thông thường mỗi đêm \\
\midrule
\multicolumn{4}{l}{\textbf{Nhóm $F_4$ (34 đặc trưng): $F_3$ + 5 nhân khẩu học \& kinh tế xã hội}} \\
RIDAGEYR & Nhân khẩu ($F_4$) & Liên tục (tuổi) & Tuổi của người tham gia khảo sát tại thời điểm khám \\
RIAGENDR & Nhân khẩu ($F_4$) & Phân loại nhị phân & Giới tính sinh học (Nam / Nữ) \\
RIDRETH1 & Nhân khẩu ($F_4$) & Phân loại 5 nhóm & Chủng tộc / Dân tộc tự xác định \\
DMDEDUC2 & Nhân khẩu ($F_4$) & Thứ bậc (1--5) & Trình độ học vấn cao nhất đạt được (người lớn) \\
INDFMPIR & Kinh tế ($F_4$) & Liên tục (tỷ số) & Tỷ số thu nhập gia đình trên chuẩn nghèo liên bang \\
\end{longtable}
\endgroup


\section{Cấu hình siêu tham số mô hình}
\label{sec:app-hyperparameters}

Mô hình cơ sở Random Forest sử dụng cấu hình siêu tham số xác định từ quy trình Nested Cross-Validation:
\begin{itemize}
    \item Số lượng cây quyết định (\textnormal{n\_estimators}): 150 cây.
    \item Độ sâu tối đa của cây (\textnormal{max\_depth}): 8 tầng.
    \item Tỷ lệ đặc trưng lấy mẫu ngẫu nhiên (\textnormal{max\_features}): 0.5 (50\% số đặc trưng).
    \item Số lượng mẫu tối thiểu để phân nhánh (\textnormal{min\_samples\_split}): 10 mẫu; tại nút lá (\textnormal{min\_samples\_leaf}): 8 mẫu.
    \item Cơ chế trọng số lớp: \textnormal{'balanced\_subsample'}; Tiêu chuẩn chia nhánh: chỉ số Gini.
\end{itemize}

Cấu hình các mô hình cơ sở khác:
\begin{itemize}
    \item \textit{Logistic Regression (LR):} Chuẩn hóa $L_2$, thuật toán tối ưu lbfgs, tham số điều hòa $C = 1.0$, trọng số lớp balanced.
    \item \textit{XGBoost (XGB):} \textnormal{n\_estimators = 100, max\_depth = 4, learning\_rate = 0.05, subsample = 0.8, colsample\_bytree = 0.8}.
    \item \textit{Ordinal Logistic Regression (OLR):} Tối ưu hóa hàm hợp lý logarit qua SciPy minimize với thuật toán L-BFGS-B, khởi tạo các ngưỡng cắt đơn điệu bằng phân vị mẫu.
\end{itemize}

\section{Kết quả chi tiết theo Fold của kiểm định chéo lồng nhau}
\label{sec:app-fold-results}

Bảng \ref{tab:app-fold-results} tổng hợp Macro-F1 trên từng fold kiểm tra ngoài mẫu (Outer Folds) của 4 thuật toán trên bài toán 5 mức $T_5$ ($F_4$).

\begin{table}[htbp]{Chi tiết hiệu năng Macro-F1 theo từng Outer Fold trên không gian $F_4$ ($T_5$).}
\centering
\footnotesize
\setlength{\tabcolsep}{4pt}
\label{tab:app-fold-results}
\begin{tabularx}{\textwidth}{l *{6}{>{\centering\arraybackslash}X}}
\toprule
\textbf{Mô hình} & \textbf{Fold 1} & \textbf{Fold 2} & \textbf{Fold 3} & \textbf{Fold 4} & \textbf{Fold 5} & \textbf{Trung bình $\pm$ SD} \\
\midrule
Logistic Regression & 0.2599 & 0.2198 & 0.2566 & 0.2382 & 0.2292 & 0.2407 $\pm$ 0.0173 \\
Decision Tree       & 0.2538 & 0.2404 & 0.2634 & 0.2239 & 0.2245 & 0.2412 $\pm$ 0.0175 \\
XGBoost (Tuned)     & 0.3017 & 0.2207 & 0.2925 & 0.2341 & 0.2615 & 0.2621 $\pm$ 0.0353 \\
Random Forest (Frozen) & 0.3111 & 0.2615 & 0.2936 & 0.2425 & 0.2690 & 0.2755 $\pm$ 0.0270 \\
\bottomrule
\end{tabularx}
\end{table}

\section{Bốn ma trận chi phí Bayes và giải thuật hiệu chỉnh đơn điệu}
\label{sec:app-cost-matrices}

Nghiên cứu khảo sát 4 ma trận chi phí bất đối xứng $\mathbf{C} \in \mathbb{R}^{3 \times 3}$:
\begin{itemize}
    \item \textit{Ma trận đồng nhất $M_0$:} $\mathbf{C}(i, j) = 1.0$ với $i \ne j$, và $\mathbf{C}(i, i) = 0$.
    \item \textit{Ma trận lâm sàng $M_1$ (Phạt vừa phải):} $\mathbf{C}(2, 0) = 3.0$; $\mathbf{C}(2, 1) = 1.5$; $\mathbf{C}(1, 0) = 1.5$; $\mathbf{C}(0, 2) = 1.5$; các lỗi khác giữ 1.0.
    \item \textit{Ma trận lâm sàng tăng cường $M_2$:} $\mathbf{C}(2, 0) = 5.0$; $\mathbf{C}(2, 1) = 2.0$; $\mathbf{C}(1, 0) = 2.0$; $\mathbf{C}(0, 2) = 2.0$.
    \item \textit{Ma trận lâm sàng tối đa $M_3$:} $\mathbf{C}(2, 0) = 8.0$; $\mathbf{C}(2, 1) = 3.0$; $\mathbf{C}(1, 0) = 2.5$; $\mathbf{C}(0, 2) = 2.5$.
\end{itemize}

Trong giải thuật phân loại nhị phân tích lũy ($E_{7b}$), hai mô hình độc lập dự đoán $\hat{P}_1 = P(Y \ge 1 \mid \mathbf{x})$ và $\hat{P}_2 = P(Y \ge 2 \mid \mathbf{x})$. Khi xuất hiện vi phạm đơn điệu ($\hat{P}_2 > \hat{P}_1$), giải thuật Isotonic Probability Clipping gán $\tilde{P}_2 = \min(\hat{P}_1, \hat{P}_2)$. Điều này bảo đảm $P(C_1 \mid \mathbf{x}) = \hat{P}_1 - \tilde{P}_2 \ge 0$, nhưng tại điểm vi phạm, $P(C_1) = 0$, giải thích hiện tượng suy giảm độ nhạy của lớp trung gian $C_1$ trên thực nghiệm.
